    \addcontentsline{toc}{section}{Appendix}
	\large
		\begingroup
			\let\clearpage\relax
			\chapter*{Appendix}
		\endgroup
    \normalsize
		\appendix

\subsection*{Appendix A: Additional System Screenshots}

The following figures provide additional screenshots of the StudyBug interface beyond those presented in Chapter 8.

\begin{figure}[H]
    \centering
    \includegraphics[width=0.9\textwidth]{img/imagin.png}
    \caption{StudyBug Dashboard Showing XP, Streaks, and Learning Progress}
\end{figure}

\begin{figure}[H]
    \centering
    \includegraphics[width=0.9\textwidth]{img/studyr.png}
    \caption{Virtual Study Room with Interactive Hotspots}
\end{figure}

\pagebreak

\subsection*{Appendix B: Database Table Schemas}

The following tables define the structure of each database table used in StudyBug. All tables include a \texttt{user\_id} foreign key for data isolation.

\begin{table}[H]
\centering
\begin{tabular}{|l|l|l|}
\hline
\textbf{Column} & \textbf{Type} & \textbf{Description} \\
\hline
id & UUID (PK) & Unique user identifier (from Supabase Auth) \\
email & VARCHAR & User email address \\
display\_name & VARCHAR & User display name \\
xp & INTEGER & Accumulated experience points \\
streak & INTEGER & Current daily streak count \\
last\_active & TIMESTAMP & Last login/activity timestamp \\
productivity\_score & FLOAT & Calculated productivity metric \\
created\_at & TIMESTAMP & Account creation timestamp \\
\hline
\end{tabular}
\caption{users Table Schema}
\end{table}

\begin{table}[H]
\centering
\begin{tabular}{|l|l|l|}
\hline
\textbf{Column} & \textbf{Type} & \textbf{Description} \\
\hline
task\_id & SERIAL (PK) & Unique task identifier \\
user\_id & UUID (FK) & Reference to user \\
title & VARCHAR & Task title \\
description & TEXT & Task description \\
priority & VARCHAR & Priority level (Low/Medium/High) \\
status & VARCHAR & Task status (pending/completed) \\
target\_date & DATE & Planned completion date \\
estimated\_pomodoros & INTEGER & Estimated Pomodoro sessions \\
completed\_pomodoros & INTEGER & Actual Pomodoro sessions completed \\
progress\_percent & INTEGER & Completion percentage (0--100) \\
created\_at & TIMESTAMP & Task creation timestamp \\
\hline
\end{tabular}
\caption{tasks Table Schema}
\end{table}

\begin{table}[H]
\centering
\begin{tabular}{|l|l|l|}
\hline
\textbf{Column} & \textbf{Type} & \textbf{Description} \\
\hline
session\_id & SERIAL (PK) & Unique session identifier \\
user\_id & UUID (FK) & Reference to user \\
task\_id & INTEGER (FK) & Associated task \\
duration\_minutes & INTEGER & Planned session duration \\
actual\_duration & INTEGER & Actual time spent (seconds) \\
status & VARCHAR & Session state (IN\_PROGRESS/COMPLETED/CANCELLED) \\
started\_at & TIMESTAMP & Session start time \\
ended\_at & TIMESTAMP & Session end time \\
\hline
\end{tabular}
\caption{focus\_sessions Table Schema}
\end{table}

\begin{table}[H]
\centering
\begin{tabular}{|l|l|l|}
\hline
\textbf{Column} & \textbf{Type} & \textbf{Description} \\
\hline
note\_id & SERIAL (PK) & Unique note identifier \\
user\_id & UUID (FK) & Reference to user \\
title & VARCHAR & Note title \\
content & TEXT & Note content \\
images & TEXT[] & Array of Base64-encoded image data \\
color & VARCHAR & Pastel color code for UI display \\
created\_at & TIMESTAMP & Note creation timestamp \\
updated\_at & TIMESTAMP & Last update timestamp \\
\hline
\end{tabular}
\caption{notes Table Schema}
\end{table}

\begin{table}[H]
\centering
\begin{tabular}{|l|l|l|}
\hline
\textbf{Column} & \textbf{Type} & \textbf{Description} \\
\hline
entry\_id & SERIAL (PK) & Unique journal entry identifier \\
user\_id & UUID (FK) & Reference to user \\
entry\_date & DATE & Date of the journal entry (unique per user) \\
content & TEXT & Journal text content \\
mood & VARCHAR & Selected mood (Happy/Sad/Tired/Excited, etc.) \\
word\_count & INTEGER & Auto-calculated word count \\
created\_at & TIMESTAMP & Entry creation timestamp \\
\hline
\end{tabular}
\caption{journal\_entries Table Schema}
\end{table}

\pagebreak

\subsection*{Appendix C: Complete API Endpoint Reference}

The following table lists all API endpoints implemented in the StudyBug backend. All routes (except authentication) require a valid JWT token in the \texttt{Authorization} header.

\begin{longtable}{|l|l|p{6cm}|}
\hline
\textbf{Method} & \textbf{Endpoint} & \textbf{Description} \\
\hline
\endfirsthead
\hline
\textbf{Method} & \textbf{Endpoint} & \textbf{Description} \\
\hline
\endhead
\multicolumn{3}{r}{\textit{Continued on next page}} \\
\endfoot
\hline
\endlastfoot
POST & /api/auth/verify & Verify Supabase JWT token \\
GET & /api/auth/profile & Get authenticated user profile \\
PUT & /api/auth/profile & Update user display name \\
\hline
GET & /api/tasks & Get all tasks for user \\
POST & /api/tasks & Create a new task \\
PUT & /api/tasks/:id & Update task details \\
PATCH & /api/tasks/:id/complete & Mark task as complete \\
DELETE & /api/tasks/:id & Soft-delete task to trash \\
\hline
GET & /api/notes & Get all notes for user \\
POST & /api/notes & Create a new note \\
PUT & /api/notes/:id & Update note content \\
DELETE & /api/notes/:id & Soft-delete note to trash \\
\hline
GET & /api/journal & Get all journal entries for user \\
POST & /api/journal & Create a journal entry \\
PUT & /api/journal/:id & Update journal entry \\
DELETE & /api/journal/:id & Delete journal entry \\
\hline
GET & /api/calendar & Get all calendar events \\
POST & /api/calendar & Create a calendar event \\
PUT & /api/calendar/:id & Update calendar event \\
DELETE & /api/calendar/:id & Delete calendar event \\
\hline
GET & /api/files & Get all uploaded files \\
POST & /api/files/upload & Upload a PDF or TXT file \\
DELETE & /api/files/:id & Delete a file \\
\hline
GET & /api/focus-sessions & Get focus session history \\
POST & /api/focus-sessions & Start a new focus session \\
PUT & /api/focus-sessions/:id & Update session status \\
GET & /api/focus-sessions/timeline & Get timeline analytics data \\
\hline
GET & /api/progress & Get user learning progress \\
PUT & /api/progress/xp & Add XP points \\
GET & /api/progress/levels & Get all learning levels \\
PUT & /api/progress/streak & Update daily streak \\
GET & /api/progress/productivity & Get productivity score \\
\hline
GET & /api/quiz/default & Get default quiz questions \\
POST & /api/quiz/generate & Generate MCQs from uploaded file \\
POST & /api/quiz/submit & Submit quiz answers and award XP \\
\hline
GET & /api/trash & Get all soft-deleted items \\
POST & /api/trash/restore/:id & Restore item from trash \\
DELETE & /api/trash/:id & Permanently delete item \\
DELETE & /api/trash/empty & Empty all trash items \\
\hline
\end{longtable}

\pagebreak

\subsection*{Appendix D: Gamification System — Level Structure}

StudyBug implements a three-phase learning progression system with 15 levels to maintain long-term user engagement.

\begin{table}[H]
\centering
\begin{tabular}{|l|c|l|c|}
\hline
\textbf{Phase} & \textbf{Level} & \textbf{Level Name} & \textbf{XP Required} \\
\hline
Phase 1: Foundations \& Basics & 1  & Curious Beginner   & 0 \\
                               & 2  & Eager Learner      & 100 \\
                               & 3  & Knowledge Seeker   & 250 \\
                               & 4  & Focused Student    & 450 \\
                               & 5  & Steady Scholar     & 700 \\
\hline
Phase 2: Building Momentum     & 6  & Rising Academic    & 1000 \\
                               & 7  & Determined Thinker & 1400 \\
                               & 8  & Skilled Analyst    & 1900 \\
                               & 9  & Sharp Mind         & 2500 \\
                               & 10 & Deep Diver         & 3200 \\
\hline
Phase 3: Peak Mastery          & 11 & Expert Learner     & 4000 \\
                               & 12 & Master Scholar     & 5000 \\
                               & 13 & Grand Thinker      & 6500 \\
                               & 14 & Elite Academic     & 8500 \\
                               & 15 & Study Legend       & 11000 \\
\hline
\end{tabular}
\caption{StudyBug Gamification Level Structure}
\end{table}

XP is awarded through the following actions:
\begin{itemize}
    \item Completing a quiz: \textbf{+50 XP} per quiz
    \item Completing a task: \textbf{+20 XP} per task
    \item Completing a focus session: \textbf{+10 XP} per session
    \item Daily login streak bonus: \textbf{+5 XP} per day
\end{itemize}

\subsection*{Appendix E: Team Contributions}

The following table summarizes the primary responsibilities of each team member during the development of StudyBug.

\begin{table}[H]
\centering
\begin{tabular}{|l|l|p{7cm}|}
\hline
\textbf{Member} & \textbf{Roll No.} & \textbf{Primary Contributions} \\
\hline
Anjana Silinchhe Shrestha & KCE080BCT004 & Frontend UI development, StudyRoom interactive hotspots, Calendar module, overall UI/UX design \\
\hline
Pradipta Joshi & KCE080BCT021 & Backend API development, authentication integration, database schema design \\
\hline
Prasant Rai & KCE080BCT024 & Pomodoro timer implementation, focus session tracking, export functionality (CSV/PDF) \\
\hline
Sushma Shrestha & KCE080BCT043 & Quiz generation system, notes and journal modules, gamification (XP, streaks, levels) \\
\hline
\end{tabular}
\caption{Team Member Contributions}
\end{table}
